\SourcePage{187}{192}
\section{Dipole radiation}

We apply the formulas obtained above to photon emission by a relativistic electron in a prescribed external field. In this case, the transition current is the matrix element of the operator
\[
\hat j=\widehat{\bar\psi}\gamma\widehat\psi,
\]
\SourcePage{188}{193}
in which the $\psi$ operators are understood to be expanded in a system of wave functions for the stationary electron states in the given field (see \S~32). The matrix element $\langle 0_i1_f|\hat j|1_i0_f\rangle$ corresponds to the electron transition from state $i$ to state $f$. This change in the occupation numbers is effected by the operator $\hat a_f^+\hat a_i$, and the transition current is therefore
\begin{equation}
j_{fi}^{\mu}=\bar\psi_f\gamma^\mu\psi_i
=(\psi_f^*\psi_i,\,\psi_f^*\boldsymbol\alpha\psi_i),
\tag{45.1}
\end{equation}
where $\psi_i$ and $\psi_f$ are the wave functions of the initial and final electron states.

Choose the photon wave function in the three-dimensionally transverse gauge (the polarization 4-vector is $e=(0,\mathbf e)$). The product in (43.10) is then $j_{fi}e^*=-\mathbf j_{fi}\mathbf e^*$. Substituting $V_{fi}$ in (44.4), we obtain the following formula for the emission probability (per second) into the solid-angle element $do$ for a photon with polarization $\mathbf e$:
\begin{equation}
dw_{\mathbf e\mathbf n}=e^2\frac\omega{2\pi}
\left|\mathbf e^*\mathbf j_{fi}(\mathbf k)\right|^2do,
\tag{45.2}
\end{equation}
where
\begin{equation}
\mathbf j_{fi}(\mathbf k)=\int\psi_f^*\boldsymbol\alpha\psi_i
e^{-i\mathbf k\mathbf r}\,d^3x.
\tag{45.3}
\end{equation}

Summation over the photon polarizations is carried out by averaging over the directions of $\mathbf e$ (in the plane perpendicular to the specified direction $\mathbf n=\mathbf k/\omega$), and then multiplying the result by 2 for the two independent possibilities for transverse photon polarization\footnote{The averaging uses the formula
\begin{equation}
\overline{e_i e_k^*}=\frac12(\delta_{ik}-n_i n_k)
\tag{45.4\text{\fontencoding{T2A}\selectfont а}}
\end{equation}
or
\begin{equation}
\overline{(\mathbf a\mathbf e)(\mathbf b\mathbf e^*)}
=\frac12\{\mathbf a\mathbf b-(\mathbf a\mathbf n)(\mathbf b\mathbf n)\}
=\frac12(\mathbf a\times\mathbf n)(\mathbf b\times\mathbf n),
\tag{45.4\text{\fontencoding{T2A}\selectfont б}}
\end{equation}
where $\mathbf a$ and $\mathbf b$ are constant vectors.}. This gives
\begin{equation}
dw_{\mathbf n}=e^2\frac\omega{2\pi}
\left|\mathbf n\times\mathbf j_{fi}(\mathbf k)\right|^2do.
\tag{45.4}
\end{equation}

Of particular importance is the case in which the photon wavelength $\lambda$ is large compared with the dimensions $a$ of the radiating system. This situation is usually connected with particle velocities small compared with the speed of light. To the first approximation in $a/\lambda$ (corresponding to dipole radiation; cf. II, \S~67), the factor $e^{-i\mathbf k\mathbf r}$ in the transition current (45.3), which varies little in the
\SourcePage{189}{194}
region where either $\psi_i$ or $\psi_f$ differs appreciably from zero, may be replaced by unity. In other words, this replacement means neglecting the photon momentum in comparison with the particle momenta in the system.

In the same approximation, the integral $\mathbf j_{fi}(0)$ may be replaced by its nonrelativistic expression, namely the matrix element $\mathbf v_{fi}$ of the electron velocity evaluated with Schrödinger wave functions. In turn, this element is $\mathbf v_{fi}=-i\omega\mathbf r_{fi}$, while $e\mathbf r_{fi}=\mathbf d_{fi}$, where $\mathbf d$ is the dipole moment of the electron (in its orbital motion). We consequently obtain the following formula for the dipole-radiation probability:
\begin{equation}
dw_{\mathbf e\mathbf n}=\frac{\omega^3}{2\pi}
\left|\mathbf e^*\mathbf d_{fi}\right|^2do
\tag{45.5}
\end{equation}
(the direction $\mathbf n$ enters here implicitly: the vector $\mathbf e$ must be perpendicular to $\mathbf n$). Summing over polarizations gives
\begin{equation}
dw_{\mathbf n}=\frac{\omega^3}{2\pi}
\left|\mathbf n\times\mathbf d_{fi}\right|^2do.
\tag{45.6}
\end{equation}

Because these formulas are nonrelativistic with respect to the electron, their generalization to arbitrary electron systems is immediate: $\mathbf d_{fi}$ is to mean the matrix element of the total dipole moment of the system.

Integrating formula (45.6) over all directions, we find the total emission probability:
\begin{equation}
w=\frac{4\omega^3}{3}|\mathbf d_{fi}|^2,
\tag{45.7}
\end{equation}
or, in ordinary units,
\begin{equation}
w=\frac{4\omega^3}{3\hbar c^3}|\mathbf d_{fi}|^2.
\tag{45.7\text{\fontencoding{T2A}\selectfont а}}
\end{equation}

The radiation intensity $I$ is obtained by multiplying the probability by $\hbar\omega$:
\begin{equation}
I=\frac{4\omega^4}{3c^3}|\mathbf d_{fi}|^2.
\tag{45.8}
\end{equation}

This formula displays a direct analogy with the classical formula (see II, (67.11)) for the intensity of dipole radiation from a system of periodically moving particles: the radiation intensity at frequency $\omega_s=s\omega$ (where $\omega$ is the frequency of the particle motion and $s$ is an integer) is
\begin{equation}
I_s=\frac{4\omega_s^4}{3c^3}|\mathbf d_s|^2,
\tag{45.9}
\end{equation}
\SourcePage{190}{195}
where $\mathbf d_s$ are the Fourier components of the dipole moment, that is, the coefficients in the expansion
\begin{equation}
\mathbf d(t)=\sum_{s=-\infty}^{\infty}\mathbf d_s e^{-is\omega t}.
\tag{45.10}
\end{equation}

The quantum formula (45.8) follows from (45.9) when these Fourier components are replaced by the matrix elements for the corresponding transitions. This rule, which expresses Bohr's \emph{correspondence principle}, is a special case of the general correspondence between Fourier components of classical quantities and quantum matrix elements in the semiclassical case (see III, \S~48). Radiation is semiclassical for transitions between states having large quantum numbers; the transition frequency $\hbar\omega=E_i-E_f$ is then small compared with the energies $E_i$ and $E_f$ of the radiating system. This circumstance would not, however, change the form of (45.8), which is valid for arbitrary transitions. It explains the fact---accidental in a certain sense---that the correspondence principle for the radiation intensity is valid not only in the semiclassical case but also in the general quantum case.
